\documentclass[12pt, titlepage]{article}

\usepackage{booktabs}
\usepackage{tabularx}
\usepackage{hyperref}
\hypersetup{
    colorlinks,
    citecolor=black,
    filecolor=black,
    linkcolor=red,
    urlcolor=blue
}
\usepackage[round]{natbib}

\input{../Comments}
\input{../Common}

\begin{document}

\title{Verification and Validation Report: \progname} 
\author{\authname}
\date{\today}
	
\maketitle

\pagenumbering{roman}

\section{Revision History}

\begin{tabularx}{\textwidth}{p{3cm}p{2cm}X}
\toprule {\bf Date} & {\bf Version} & {\bf Notes}\\
\midrule
Date 1 & 1.0 & Notes\\
Date 2 & 1.1 & Notes\\
\bottomrule
\end{tabularx}

~\newpage

\section{Symbols, Abbreviations and Acronyms}

\renewcommand{\arraystretch}{1.2}
\begin{tabular}{l l} 
  \toprule		
  \textbf{symbol} & \textbf{description}\\
  \midrule 
  T & Test\\
  \bottomrule
\end{tabular}\\

\wss{symbols, abbreviations or acronyms -- you can reference the SRS tables if needed}

\newpage

\tableofcontents

\listoftables %if appropriate

\listoffigures %if appropriate

\newpage

\pagenumbering{arabic}

This document ...

\section{Functional Requirements Evaluation}

\section{Nonfunctional Requirements Evaluation}

\subsection{Usability}
		
\subsection{Performance}

\subsection{etc.}
	
\section{Comparison to Existing Implementation}

To better understand the effectiveness of the proposed system,
Hitchly, it is useful to compare it with existing ride-sharing
platforms such as Uber, Lyft, and Poparide. These systems provide
transportation services using mobile applications and large driver
networks, ranging from carpooling to on-demand rides. While the core
concept of connecting riders and drivers is similar, the goals and
design of Hitchly differ in several important ways.

\subsection{Target Users}

Uber, Lyft, and Poparide are designed for the general public and
operate on a global scale. Hitchly is specifically designed for
university students and faculty at McMaster University, focusing on
a smaller, community-based user group. Restricting the system to
verified McMaster University users increases trust between riders
and drivers.

\subsection{Matching Model}

Uber, Lyft, and Poparide share a similarity in that they do not
emphasize driver–rider interaction beyond the ride itself. Hitchly
not only connects riders and drivers but also allows users to specify
preferences related to personality or behavior. Uber and Lyft also
differ from Hitchly because they use real-time driver dispatch
systems where riders are matched with the nearest available driver.
Poparide and Hitchly share a similar approach in that users have the
ability to offer or request rides that align with planned trips.

\subsection{Cost Structure}

Commercial ride-sharing platforms primarily focus on generating
profit for drivers and the platform itself. In contrast, Hitchly is
intended to reduce transportation costs for students through ride
sharing and cost splitting.

\subsection{Summary}

While Uber and Lyft provide large-scale on-demand transportation
services, Hitchly focuses on community-based ride coordination for
students. By emphasizing route matching, scheduled trips, and cost
sharing, the system provides a solution better suited for university
commuting needs.


\section{Unit Testing}

\section{Changes Due to Testing}

\wss{This section should highlight how feedback from the users and from 
the supervisor (when one exists) shaped the final product.  In particular 
the feedback from the Rev 0 demo to the supervisor (or to potential users) 
should be highlighted.}

\section{Automated Testing}
		
\section{Trace to Requirements}
		
\section{Trace to Modules}		

\section{Code Coverage Metrics}

\bibliographystyle{plainnat}
\bibliography{../../refs/References}

\newpage{}
\section*{Appendix --- Reflection}

The information in this section will be used to evaluate the team members on the
graduate attribute of Reflection.

\input{../Reflection.tex}

\begin{enumerate}
  \item What went well while writing this deliverable? 
  \item What pain points did you experience during this deliverable, and how
    did you resolve them?
  \item Which parts of this document stemmed from speaking to your client(s) or
  a proxy (e.g. your peers)? Which ones were not, and why?
  \item In what ways was the Verification and Validation (VnV) Plan different
  from the activities that were actually conducted for VnV?  If there were
  differences, what changes required the modification in the plan?  Why did
  these changes occur?  Would you be able to anticipate these changes in future
  projects?  If there weren't any differences, how was your team able to clearly
  predict a feasible amount of effort and the right tasks needed to build the
  evidence that demonstrates the required quality?  (It is expected that most
  teams will have had to deviate from their original VnV Plan.)
\end{enumerate}

\end{document}